\subsection{Three flavour neutrino oscillations in models with large extra dimensions --- arXiv:hep-ph/0006278}

\textbf{Authors:} R. N. Mohapatra, A. P\'erez-Lorenzana

\paragraph{主な主張}
3世代のactive neutrinoとbulk singletからなる最小模型は、当時の太陽・大気ニュートリノデータを説明する構成を持つ一方、LSNDを含む三者の同時説明はできないことを示す\cite{Mohapatra:2000ThreeFlavour}。

\paragraph{新規性}
3世代それぞれのKK塔との弱混合・強混合を分類し、小さなニュートリノ質量を得ることと観測されたflavour転換を再現することが別の条件であると明確にした。

\paragraph{位置付け}
最小bulk neutrino模型のflavour構造と限界を整理した初期文献であり、sterileへの消失だけで短基線appearanceまで説明できるとは限らないことを示す。
